
\section{假设检验理论与似然比检验}

在计量经济学分析与时间序列建模中，对模型参数稳定性的统计推断具有基础性意义。无论是检验宏观经济系统的结构是否发生变化，还是识别金融变量之间动态关系是否出现断裂，均需要建立在明确的假设检验框架之上。

本节从经典统计推断中的假设检验原理出发，进一步引入似然比检验的基本理论，并讨论其在向量自回归模型已知断点检验中的应用。随后，针对高维条件下经典渐近理论失效的问题，进一步介绍基于 Bootstrap 的非渐近重抽样推断思路。

\subsection{假设检验基本原理}

假设检验的基本任务，是利用样本所提供的有限信息，对总体分布或参数所满足的某一命题作出统计决策。设概率空间为 $(\Omega,\mathcal{F},P)$，给定样本
\begin{equation}
X=(X_1,X_2,\dots,X_n),
\end{equation}
其联合密度函数或概率质量函数记为 $f(x;\theta)$，其中参数 $\theta$ 属于参数空间 $\Theta$。在统计推断中，通常将参数空间划分为两个互不相交的子集 $\Theta_0$ 与 $\Theta_1$，满足
\begin{equation}
\Theta_0\cup\Theta_1=\Theta, \qquad \Theta_0\cap\Theta_1=\emptyset.
\end{equation}
据此，原假设与备择假设分别写为
\begin{equation}
H_0: \theta\in\Theta_0, \qquad H_1: \theta\in\Theta_1.
\end{equation}

假设检验的逻辑基础是\textbf{小概率原理}与\textbf{反证法}的结合。其推理过程如下：首先假定原假设 $H_0$ 为真，在此前提下推导检验统计量 $T(X)$ 的抽样分布；然后观察实际样本所计算出的统计量值是否落入该分布的极端尾部区域。若统计量取值在 $H_0$ 成立的条件下属于概率极低的事件（即落入预先设定的拒绝域 $R$），则依据小概率事件在单次试验中不应发生的原理，认为 $H_0$ 的假定与观测数据相矛盾，从而拒绝原假设；反之，若统计量未落入拒绝域，则表明当前样本与 $H_0$ 并不矛盾，不拒绝原假设。需要强调的是，”不拒绝”并不等同于原假设为真，而仅表示在当前样本信息下尚无充分的统计证据反对 $H_0$。

\subsection{假设检验核心评价指标}

由于样本具有随机性，假设检验不可避免地会产生两类错误。

第一类错误是指当原假设成立时错误地拒绝原假设，其概率记为
\begin{equation}
\alpha = \sup_{\theta\in\Theta_0} P(X\in R\mid \theta),
\end{equation}
通常称为检验的显著性水平或规模（size）。这里取上确界而非普通期望，是因为 $\Theta_0$ 可能包含多个参数值，不同参数下第一类错误率各不相同；取上确界保证了无论真实参数在 $\Theta_0$ 中取何值，检验犯第一类错误的概率均不超过 $\alpha$，从而为检验的虚警控制提供了最坏情形保障。

第二类错误是指当备择假设成立时未能拒绝原假设，其概率记为
\begin{equation}
\beta = P(X\notin R\mid \theta\in\Theta_1).
\end{equation}
与之对应，统计功效（power）定义为在备择假设成立时正确拒绝原假设的概率，即
\begin{equation}
1-\beta = P(X\in R\mid \theta\in\Theta_1).
\end{equation}
功效是衡量检验灵敏度的核心指标：功效越高，检验越能在真实参数确实偏离原假设时将其识别出来；反之，功效不足的检验即使面对实质性的参数变化也倾向于"漏检"，从而给出过于保守的结论。

在有限样本条件下，第一类错误与第二类错误之间通常存在权衡关系。若为了降低第一类错误而过度收缩拒绝域，则会提高第二类错误，进而削弱检验功效。因此，经典假设检验的基本目标是在给定显著性水平 $\alpha$ 的前提下，尽可能提高检验功效。

\subsection{似然比检验理论及其高维局限性}

在大多数实际计量分析中，参数空间往往包含多个未知参数，因此检验问题通常表现为复合假设检验。为处理此类情形，统计学中引入了广义似然比检验（Generalized Likelihood Ratio Test, GLRT）。其基本思想是，分别在受约束参数空间与不受约束参数空间中对似然函数进行最大化，并通过比较两者的最大似然值，评估原假设所施加约束对模型拟合优度的影响。

设样本似然函数为 $L(\theta;x)$，则广义似然比统计量定义为
\begin{equation}
\Lambda(x)=\frac{\sup_{\theta\in\Theta_0}L(\theta;x)}{\sup_{\theta\in\Theta}L(\theta;x)}.
\end{equation}
由于 $\Theta_0\subseteq\Theta$，故有 $0\leq \Lambda(x)\leq 1$。当 $\Lambda(x)$ 接近于 $1$ 时，表明原假设约束下的模型拟合与全模型差异较小，样本数据不足以反驳原假设；反之，当 $\Lambda(x)$ 接近于 $0$ 时，则说明原假设显著损害了模型对数据的解释能力，应倾向于拒绝原假设。

为了便于推断，通常对 $\Lambda(x)$ 进行对数变换，并构造统计量
\begin{equation}
G=-2\ln\Lambda(x)=2\left[\ln L(\hat{\theta}_{\mathrm{MLE}};x)-\ln L(\hat{\theta}_0;x)\right],
\end{equation}
其中，$\hat{\theta}_{\mathrm{MLE}}$ 为无约束极大似然估计，$\hat{\theta}_0$ 为原假设约束下的极大似然估计。根据威尔克斯定理（Wilks' theorem），在样本量趋于无穷且满足一系列正则性条件的低维环境下，统计量 $G$ 的极限分布为卡方分布：
\begin{equation}
G \xrightarrow{d} \chi_r^2,
\end{equation}
其中，$r$ 为受约束模型相对于全模型所施加的独立约束个数。

在时间序列分析中，似然比检验常被用于结构变化识别。以向量自回归模型的已知断点检验为例，考虑一个包含 $N$ 个内生变量、滞后阶数为 $p$ 的 VAR($p$) 模型：
\begin{equation}
Y_t=c+\Phi_1Y_{t-1}+\cdots+\Phi_pY_{t-p}+\epsilon_t,
\end{equation}
其中，误差项满足多元正态分布，协方差矩阵为 $\Sigma$。

\paragraph{假设设定（$H_0$ 与 $H_1$）}若已知样本中的某一时间点 $\tau$ 可能对应结构变化，则原假设 $H_0$ 表示整个样本期内参数保持稳定；备择假设 $H_1$ 表示系统在时间点 $\tau$ 发生结构性断裂，断点前后分别对应不同的参数矩阵。

\paragraph{统计量构造}在多元正态假设下，忽略常数项后，对数似然函数的核心部分可表示为残差协方差矩阵行列式的函数。对于长度为 $T$ 的样本，其极大对数似然可写为
\begin{equation}
\ln L \propto -\frac{T}{2}\ln|\hat{\Sigma}|.
\end{equation}
因此，在原假设下，利用全样本拟合受约束模型可得到残差协方差矩阵 $\hat{\Sigma}_R$；在备择假设下，对断点前后两个子样本分别拟合无约束模型，可得到残差协方差矩阵 $\hat{\Sigma}_1$ 与 $\hat{\Sigma}_2$。由此可构造已知断点下的对数似然比统计量：
\begin{equation}
LR = T_{\mathrm{eff}}\ln|\hat{\Sigma}_R| - T_1\ln|\hat{\Sigma}_1| - T_2\ln|\hat{\Sigma}_2|,
\end{equation}
其中，$T_{\mathrm{eff}}$ 为有效样本长度，$T_1$ 与 $T_2$ 分别表示断点前后两个子样本的长度。

\paragraph{渐近分布}在低维设定下，若 $N$ 与 $p$ 均较小，且样本规模足够大，则该统计量可近似服从卡方分布，其自由度等于断裂前后允许变化的参数个数。对于允许全部系数发生变化的 VAR($p$) 模型，约束数一般为
\begin{equation}
k = N(Np+1)=N^2p+N.
\end{equation}

\paragraph{决策规则}在经典低维框架中，研究者依据 $\chi_k^2$ 分布确定临界值：若 $LR$ 超过显著性水平 $\alpha$ 对应的临界值，则拒绝参数稳定的原假设，认定系统在 $\tau$ 处发生了结构性变化；否则不拒绝 $H_0$。

然而，上述渐近卡方分布的成立依赖于一个关键前提：参数维度相对于样本量而言足够小。在高维 VAR 模型中，参数总数以 $N^2p$ 的速度增长，当变量维度 $N$ 较大而样本长度 $T$ 有限时，待估参数数量与样本量之比不再接近于零，威尔克斯定理的近似精度随之急剧恶化。直观地说，参数过多导致非约束模型有"余力"去拟合样本中的随机噪声，使得两个模型的似然值差距被人为放大，似然比统计量因此系统性地偏大，不再服从其名义上对应的卡方分布。

高维条件下，若进一步引入 Lasso 稀疏惩罚或低秩约束来估计 VAR 模型，渐近推断的困难则更为根本。Lasso 的 $L_1$ 惩罚项在零点不可微，使得估计量的抽样分布高度不规则；低秩约束通过奇异值截断强制将参数映射到低维流形，同样破坏了参数空间的凸性与光滑性。在这两种正则化约束下，检验统计量不再收敛于任何具有闭式表达的标准分布，卡方临界值表完全失去适用性。因此，在高维正则化 VAR 模型中，必须寻求不依赖渐近分布的推断方法。

\subsection{Bootstrap 自助法重抽样推断理论}

Bootstrap 的核心思想是将统计量的抽样分布问题转化为一个可以通过计算模拟求解的经验问题。经典推断方法依赖于样本量趋于无穷时的渐近分布结果，而 Bootstrap 则转换思路：既然总体分布未知，不妨以样本的经验分布作为总体的近似，并在此基础上通过重复重抽样来逼近统计量的有限样本分布。

具体而言，设感兴趣的统计量为 $T_n$，其精确分布通常未知。Bootstrap 通过以下方式构造 $T_n$ 分布的经验近似：从原始样本中根据数据的依赖结构选择合适的重抽样方式生成伪样本，对每一个伪样本计算相同的统计量，经 $B$ 次重复后得到统计量的经验分布。根据格利文科定理（Glivenko--Cantelli theorem），经验分布函数依概率一致收敛于总体分布函数，当 $B$ 足够大时该经验分布即可作为真实分布的可靠近似，从而为置信区间构造、假设检验等推断任务提供数据驱动的参考分布。算法~\ref{alg:bootstrap-general} 给出了 Bootstrap 重抽样的一般流程。

\begin{algorithm}[htbp]
\caption{Bootstrap 重抽样的一般流程}\label{alg:bootstrap-general}
\KwIn{原始样本 $\mathcal{X}=\{X_1, X_2, \dots, X_n\}$，感兴趣的统计量函数 $T(\cdot)$，重抽样次数 $B$}
\KwOut{统计量 $T$ 的经验分布 $\{T^{*(1)}, T^{*(2)}, \dots, T^{*(B)}\}$}
\BlankLine
\tcc{第一步：计算原始统计量}
在原始样本 $\mathcal{X}$ 上计算观测统计量 $T_{\mathrm{obs}} \leftarrow T(\mathcal{X})$\;
\BlankLine
\tcc{第二步：重抽样循环}
\For{$b = 1, 2, \dots, B$}{
    根据数据的依赖结构选择合适的重抽样方式，从原始样本中有放回地抽取 $n$ 个观测，生成伪样本 $\mathcal{X}^{*(b)}$\;
    在伪样本 $\mathcal{X}^{*(b)}$ 上计算统计量 $T^{*(b)} \leftarrow T(\mathcal{X}^{*(b)})$\;
}
\BlankLine
\tcc{第三步：构造经验分布}
以 $\{T^{*(1)}, T^{*(2)}, \dots, T^{*(B)}\}$ 作为统计量 $T$ 的经验抽样分布，用于后续的置信区间构造、假设检验或偏差校正等推断任务\;
\end{algorithm}

